\anexo{Resultados de la encuesta}
\label{anexo:encuesta}

Se reproduce el instrumento cuantitativo aplicado y la totalidad de las tablas de frecuencia obtenidas. El análisis de los resultados y su interpretación se desarrollan en la Sección~\ref{sec:encuestas}.

\section*{Ficha técnica}
\addcontentsline{toc}{subsection}{Ficha técnica}

\begin{center}
  \begin{tabular}{ll}
    \hline
    Población objetivo & Personas de 18 a 40 años residentes en Argentina, usuarias \\
                       & de Twitter/X que siguen noticias de política o economía \\
    Técnica            & Cuestionario autoadministrado de catorce preguntas \\
    Muestreo           & No probabilístico, por conveniencia y bola de nieve \\
    Distribución       & Digital y presencial \\
    Período de campo   & 4 de julio a 18 de agosto de 2026 \\
    Respuestas obtenidas & 140 \\
    Respuestas dentro del segmento & 107 \\
    Pregunta abierta respondida & 34 de 140 \\
    \hline
  \end{tabular}
\end{center}

Las preguntas 1 a 3 se computan sobre las 140 respuestas obtenidas, por ser las que definen el filtro de segmento. A partir de la pregunta 4, la base de cálculo es de 107 casos, correspondientes a quienes declararon de forma simultánea tener entre 18 y 40 años, usar Twitter/X y seguir noticias de política o economía en la plataforma. Cada tabla indica su base.

\section*{Cuestionario aplicado}
\addcontentsline{toc}{subsection}{Cuestionario aplicado}

\textbf{Bloque A. Perfil y filtro de segmento}

\begin{enumerate}
  \item ¿Qué edad tenés? \emph{(opción única)} Menos de 18 / 18-25 / 26-33 / 34-40 / Más de 40.
  \item ¿Con qué frecuencia usás Twitter/X? \emph{(opción única)} Varias veces al día / Una vez al día / Algunas veces por semana / Rara vez / No la uso.
  \item ¿Seguís noticias de política o economía en Twitter/X? \emph{(opción única)} Sí, frecuentemente / A veces / No.
\end{enumerate}

\textbf{Bloque B. Exposición al problema}

\begin{enumerate}\setcounter{enumi}{3}
  \item ¿Con qué frecuencia sentís que te cruzás con noticias o afirmaciones falsas, engañosas o no verificables en Twitter/X? \emph{(opción única)} Muy frecuentemente / Frecuentemente / A veces / Rara vez / Nunca.
  \item ¿Alguna vez compartiste, o estuviste a punto de compartir, contenido que después resultó falso o dudoso? \emph{(opción única)} Sí, lo compartí / Estuve a punto pero no lo hice / No / No sé.
  \item ¿Qué tan preocupado/a estás por la desinformación en las redes sociales en Argentina? \emph{(escala de 1 a 5, donde 1 es nada preocupado/a y 5 es muy preocupado/a)}.
\end{enumerate}

\textbf{Bloque C. Capacidad y comportamiento actual}

\begin{enumerate}\setcounter{enumi}{6}
  \item ¿Qué tan capaz te sentís de distinguir por tu cuenta una noticia falsa de una verdadera en Twitter/X? \emph{(escala de 1 a 5, donde 1 es nada capaz y 5 es muy capaz)}.
  \item Cuando dudás de una publicación, ¿qué hacés habitualmente? \emph{(opción múltiple)} Busco en Google / Comparo con medios que conozco / Miro la cuenta que lo publicó / Le pregunto a alguien / Nada, sigo de largo / Otro.
  \item En promedio, ¿cuánto tiempo le dedicás a verificar una publicación que te genera dudas? \emph{(opción única)} Menos de 1 minuto / 1-3 minutos / Más de 3 minutos / No verifico.
\end{enumerate}

\textbf{Bloque D. Apetito y confianza en la solución}

\begin{enumerate}\setcounter{enumi}{9}
  \item Si existiera una extensión gratuita de Chrome que, al leer un tuit, te muestra un puntaje de probabilidad de que sea desinformación junto con enlaces a fuentes que lo corroboran o contradicen, ¿qué tan probable es que la instales y uses? \emph{(escala de 1 a 5, donde 1 es nada probable y 5 es muy probable)}.
  \item ¿Qué tanto confiarías en un puntaje generado automáticamente por un sistema de inteligencia artificial para esto? \emph{(escala de 1 a 5, donde 1 es no confiaría y 5 es confiaría plenamente)}.
  \item ¿Qué te haría confiar más en una herramienta así? \emph{(opción múltiple)} Que muestre las fuentes/evidencia / Que explique por qué llega a ese resultado / Que sea transparente sobre su margen de error / Que la respalde una institución reconocida / Que no tenga sesgo político / Otro.
  \item ¿Cuál sería tu mayor preocupación con una herramienta así? \emph{(opción múltiple)} Que se equivoque (falsos positivos) / Sesgo político / Privacidad de mis datos / Que censure opiniones o sátira / Ninguna / Otro.
  \item En una frase: ¿qué te resultaría más útil de una herramienta que te ayude a detectar desinformación en Twitter/X? \emph{(respuesta abierta, opcional)}.
\end{enumerate}

\section*{Tablas de frecuencia}
\addcontentsline{toc}{subsection}{Tablas de frecuencia}

\begin{table}[t]
  \centering
  \begin{tabular}{lrr}
    \hline
    \textbf{Respuesta} & \textbf{Casos} & \textbf{Porcentaje} \\
    \hline
    Menos de 18 & 3 & 2,1\% \\
    18-25 & 83 & 59,3\% \\
    26-33 & 39 & 27,9\% \\
    34-40 & 8 & 5,7\% \\
    Más de 40 & 7 & 5,0\% \\
    \hline
  \end{tabular}
  \caption{Pregunta 1. ¿Qué edad tenés? (n = 140). Fuente: elaboración propia.}
  \label{tab:enc-p1}
\end{table}

\begin{table}[t]
  \centering
  \begin{tabular}{lrr}
    \hline
    \textbf{Respuesta} & \textbf{Casos} & \textbf{Porcentaje} \\
    \hline
    Varias veces al día & 48 & 34,3\% \\
    Una vez al día & 38 & 27,1\% \\
    Algunas veces por semana & 38 & 27,1\% \\
    Rara vez & 12 & 8,6\% \\
    No la uso & 4 & 2,9\% \\
    \hline
  \end{tabular}
  \caption{Pregunta 2. ¿Con qué frecuencia usás Twitter/X? (n = 140). Fuente: elaboración propia.}
  \label{tab:enc-p2}
\end{table}

\begin{table}[t]
  \centering
  \begin{tabular}{lrr}
    \hline
    \textbf{Respuesta} & \textbf{Casos} & \textbf{Porcentaje} \\
    \hline
    Sí, frecuentemente & 54 & 38,6\% \\
    A veces & 66 & 47,1\% \\
    No & 20 & 14,3\% \\
    \hline
  \end{tabular}
  \caption{Pregunta 3. ¿Seguís noticias de política o economía en Twitter/X? (n = 140). Fuente: elaboración propia.}
  \label{tab:enc-p3}
\end{table}

\begin{table}[t]
  \centering
  \begin{tabular}{lrr}
    \hline
    \textbf{Respuesta} & \textbf{Casos} & \textbf{Porcentaje} \\
    \hline
    Muy frecuentemente & 45 & 42,1\% \\
    Frecuentemente & 34 & 31,8\% \\
    A veces & 15 & 14,0\% \\
    Rara vez & 7 & 6,5\% \\
    Nunca & 6 & 5,6\% \\
    \hline
  \end{tabular}
  \caption{Pregunta 4. ¿Con qué frecuencia sentís que te cruzás con noticias o afirmaciones falsas, engañosas o no verificables en Twitter/X? (n = 107). Fuente: elaboración propia.}
  \label{tab:enc-p4}
\end{table}

\begin{table}[t]
  \centering
  \begin{tabular}{lrr}
    \hline
    \textbf{Respuesta} & \textbf{Casos} & \textbf{Porcentaje} \\
    \hline
    Sí, lo compartí & 38 & 35,5\% \\
    Estuve a punto pero no lo hice & 28 & 26,2\% \\
    No & 28 & 26,2\% \\
    No sé & 13 & 12,1\% \\
    \hline
  \end{tabular}
  \caption{Pregunta 5. ¿Alguna vez compartiste, o estuviste a punto de compartir, contenido que después resultó falso o dudoso? (n = 107). Fuente: elaboración propia.}
  \label{tab:enc-p5}
\end{table}

\begin{table}[t]
  \centering
  \begin{tabular}{lrr}
    \hline
    \textbf{Respuesta} & \textbf{Casos} & \textbf{Porcentaje} \\
    \hline
    1 & 0 & 0,0\% \\
    2 & 4 & 3,7\% \\
    3 & 17 & 15,9\% \\
    4 & 40 & 37,4\% \\
    5 & 46 & 43,0\% \\
    \hline
  \end{tabular}
  \caption{Pregunta 6. ¿Qué tan preocupado/a estás por la desinformación en las redes sociales en Argentina? (1 = nada preocupado/a, 5 = muy preocupado/a) (n = 107). Media: 4,20. Fuente: elaboración propia.}
  \label{tab:enc-p6}
\end{table}

\begin{table}[t]
  \centering
  \begin{tabular}{lrr}
    \hline
    \textbf{Respuesta} & \textbf{Casos} & \textbf{Porcentaje} \\
    \hline
    1 & 3 & 2,8\% \\
    2 & 16 & 15,0\% \\
    3 & 28 & 26,2\% \\
    4 & 36 & 33,6\% \\
    5 & 24 & 22,4\% \\
    \hline
  \end{tabular}
  \caption{Pregunta 7. ¿Qué tan capaz te sentís de distinguir por tu cuenta una noticia falsa de una verdadera en Twitter/X? (1 = nada capaz, 5 = muy capaz) (n = 107). Media: 3,58. Fuente: elaboración propia.}
  \label{tab:enc-p7}
\end{table}

\begin{table}[t]
  \centering
  \begin{tabular}{lrr}
    \hline
    \textbf{Respuesta} & \textbf{Casos} & \textbf{Porcentaje} \\
    \hline
    Busco en Google & 89 & 83,2\% \\
    Miro la cuenta que lo publicó & 76 & 71,0\% \\
    Comparo con medios que conozco & 51 & 47,7\% \\
    Le pregunto a alguien & 45 & 42,1\% \\
    Nada, sigo de largo & 7 & 6,5\% \\
    Leo los comentarios o las notas de la comunidad si hay & 1 & 0,9\% \\
    Lo busco en X con otras palabras & 1 & 0,9\% \\
    \hline
  \end{tabular}
  \caption{Pregunta 8. Cuando dudás de una publicación, ¿qué hacés habitualmente? (opción múltiple) (n = 107). Los porcentajes se calculan sobre el total de casos y suman más de cien por tratarse de una pregunta de opción múltiple. Fuente: elaboración propia.}
  \label{tab:enc-p8}
\end{table}

\begin{table}[t]
  \centering
  \begin{tabular}{lrr}
    \hline
    \textbf{Respuesta} & \textbf{Casos} & \textbf{Porcentaje} \\
    \hline
    Menos de 1 minuto & 38 & 35,5\% \\
    1-3 minutos & 43 & 40,2\% \\
    Más de 3 minutos & 18 & 16,8\% \\
    No verifico & 8 & 7,5\% \\
    \hline
  \end{tabular}
  \caption{Pregunta 9. En promedio, ¿cuánto tiempo le dedicás a verificar una publicación que te genera dudas? (n = 107). Fuente: elaboración propia.}
  \label{tab:enc-p9}
\end{table}

\begin{table}[t]
  \centering
  \begin{tabular}{lrr}
    \hline
    \textbf{Respuesta} & \textbf{Casos} & \textbf{Porcentaje} \\
    \hline
    1 & 3 & 2,8\% \\
    2 & 7 & 6,5\% \\
    3 & 16 & 15,0\% \\
    4 & 36 & 33,6\% \\
    5 & 45 & 42,1\% \\
    \hline
  \end{tabular}
  \caption{Pregunta 10. Si existiera una extensión gratuita de Chrome que, al leer un tuit, te muestra un puntaje de probabilidad de que sea desinformación junto con enlaces a fuentes que lo corroboran o contradicen, ¿qué tan probable es que la instales y uses? (1 = nada probable, 5 = muy probable) (n = 107). Media: 4,06. Fuente: elaboración propia.}
  \label{tab:enc-p10}
\end{table}

\begin{table}[t]
  \centering
  \begin{tabular}{lrr}
    \hline
    \textbf{Respuesta} & \textbf{Casos} & \textbf{Porcentaje} \\
    \hline
    1 & 7 & 6,5\% \\
    2 & 14 & 13,1\% \\
    3 & 24 & 22,4\% \\
    4 & 35 & 32,7\% \\
    5 & 27 & 25,2\% \\
    \hline
  \end{tabular}
  \caption{Pregunta 11. ¿Qué tanto confiarías en un puntaje generado automáticamente por un sistema de inteligencia artificial para esto? (1 = no confiaría, 5 = confiaría plenamente) (n = 107). Media: 3,57. Fuente: elaboración propia.}
  \label{tab:enc-p11}
\end{table}

\begin{table}[t]
  \centering
  \begin{tabular}{lrr}
    \hline
    \textbf{Respuesta} & \textbf{Casos} & \textbf{Porcentaje} \\
    \hline
    Que no tenga sesgo político & 90 & 84,1\% \\
    Que muestre las fuentes/evidencia & 71 & 66,4\% \\
    Que explique por qué llega a ese resultado & 54 & 50,5\% \\
    Que sea transparente sobre su margen de error & 41 & 38,3\% \\
    Que la respalde una institución reconocida & 34 & 31,8\% \\
    Que sea open source & 4 & 3,7\% \\
    \hline
  \end{tabular}
  \caption{Pregunta 12. ¿Qué te haría confiar más en una herramienta así? (opción múltiple) (n = 107). Los porcentajes se calculan sobre el total de casos y suman más de cien por tratarse de una pregunta de opción múltiple. Fuente: elaboración propia.}
  \label{tab:enc-p12}
\end{table}

\begin{table}[t]
  \centering
  \begin{tabular}{lrr}
    \hline
    \textbf{Respuesta} & \textbf{Casos} & \textbf{Porcentaje} \\
    \hline
    Que se equivoque (falsos positivos) & 71 & 66,4\% \\
    Que censure opiniones o sátira & 61 & 57,0\% \\
    Sesgo político & 52 & 48,6\% \\
    Privacidad de mis datos & 38 & 35,5\% \\
    Ninguna & 3 & 2,8\% \\
    Que se vuelva paga & 3 & 2,8\% \\
    Que ande lento & 1 & 0,9\% \\
    \hline
  \end{tabular}
  \caption{Pregunta 13. ¿Cuál sería tu mayor preocupación con una herramienta así? (opción múltiple) (n = 107). Los porcentajes se calculan sobre el total de casos y suman más de cien por tratarse de una pregunta de opción múltiple. Fuente: elaboración propia.}
  \label{tab:enc-p13}
\end{table}

\clearpage

\section*{Respuestas a la pregunta abierta}
\addcontentsline{toc}{subsection}{Respuestas a la pregunta abierta}

Se transcriben las respuestas recibidas a la pregunta 14, dentro del segmento, tal como fueron consignadas.

\begin{itemize}\itemsep2pt
  \item Que te muestre posts donde se evidencia la verdad sobre la fake news
  \item Efectividad
  \item detectar imagenes generadas con ia
  \item Que muestre las dos campanas y no una sola
  \item detectar imagenes generadas con ia
  \item Efectividad y que sea rapida
  \item Que distinga entre una opinion y un dato falso
  \item Ahorrarme el tiempo de googlear cada cosa
  \item La verdad no se, capaz que me sirva para no comerme fake news
  \item no lo se
  \item Que sea transparente con como llega al resultado
  \item no lo se
  \item detectar imagenes generadas con ia
  \item Que no dependa de un solo medio para decidir si algo es falso
  \item Que me diga cuando algo es viejo y lo estan recirculando como nuevo
  \item Que no me llene la pantalla de carteles
  \item Ver el contexto que falta en el tuit
  \item Que sea facil de usar, sin tener que configurar nada
  \item La verdad no se, capaz que me sirva para no comerme fake news
  \item Que me diga cuando algo es viejo y lo estan recirculando como nuevo
  \item Que avise de titulos clickbait aunque la nota no sea falsa
  \item Que me diga cuando algo es viejo y lo estan recirculando como nuevo
  \item no lo se
  \item La verdad no se, capaz que me sirva para no comerme fake news
  \item Que avise de titulos clickbait aunque la nota no sea falsa
\end{itemize}
